% ===========================================================================
%  sviluppatore/05-estendere/03-approccio.tex
% ===========================================================================
\section{Un approccio}\label{sez:nuovo-approccio}

Un approccio \emph{scavalca il motore}. Non implementa
\simbolo{AtomicStrategy}, non emette \simbolo{Signal}, non passa per
allocatore, rischio ed esecuzione. Costruisce il proprio calcolo e consegna
direttamente la serie dei rendimenti.

\begin{center}
\begin{tabular}{@{}L{28mm}L{52mm}L{54mm}@{}}
\toprule
 & \textbf{Strategia} & \textbf{Approccio} \\
\midrule
\rigaTenue
File & \file{strategy.py} & \file{approach.py} \\
Oggetto & \cod{DESCRIPTOR} & \cod{APPROACH\_DESCRIPTOR} \\
\rigaTenue
Entrata & \cod{factory(params)} \freccia{} istanza & \cod{build(params)}
  \freccia{} \cod{pd.Series} \\
Chi esegue & il motore di backtest & l'approccio stesso \\
\rigaTenue
Produce & segnali, poi posizioni & rendimenti giornalieri \\
Runner & \file{execution/run\_walkforward.py} &
  \file{execution/run\_approach.py} \\
\bottomrule
\end{tabular}
\end{center}

\subsection{Il contratto}

\begin{codiceInterno}{python}{Dichiarato nel docstring del modulo}{approaches/descriptor.py:17}
CONTRATTO della funzione build:
    build(params: dict) -> pandas.Series
        - index = date (tz-naive); valori = rendimenti DAILY (frazione)
        - usa SOLO dati fino a t per ogni t (niente lookahead)
\end{codiceInterno}

\noindent
«Tz-naive» non è pignoleria: una serie con fuso orario e una senza non si
allineano, e il risultato è un ensemble che perde silenziosamente le
osservazioni di uno stream. Per questo gli approcci installati chiudono tutti
allo stesso modo:

\begin{codiceInterno}{python}{La chiusura obbligata di build}{approaches/momentum\_ema\_approach/approach.py:74}
    # Assicuriamo che l'indice sia tz-naive come richiesto
    pnl.index = pd.to_datetime(pnl.index).tz_localize(None)
    return pnl
\end{codiceInterno}

\subsection{Il descrittore}

Più corto di quello delle strategie, e con tre campi che si derivano da soli
se li si omette:

\begin{codiceInterno}{python}{I default derivati}{approaches/descriptor.py:79}
        if not self.label:
            self.label = self.name
        if not self.output_path_pattern:
            self.output_path_pattern = f"output/strategies/{self.name}_pnl_daily.parquet"
        if not self.backtest_script:
            self.backtest_script = f"approaches/{self.name}/run_backtest.py"
\end{codiceInterno}

\noindent
Due campi hanno invece una lista chiusa di valori ammessi, e un valore fuori
lista solleva \cod{ValueError} al caricamento:

\begin{codiceInterno}{python}{Le due allowlist}{approaches/descriptor.py:31}
APPROACH_STREAM_TYPES = ("pnl_daily", "pnl_sparse", "equity_curve")
APPROACH_STREAM_CATEGORIES = ("ml", "d1", "intraday", "ml_overlay")
\end{codiceInterno}

\noindent
Un esempio completo:

\begin{codiceInterno}{python}{Un descrittore di approccio}{approaches/momentum\_ema\_approach/approach.py:79}
APPROACH_DESCRIPTOR = ApproachDescriptor(
    name=NAME,
    label="EMA Momentum Approach",
    description=("Generates weights from the distance between the closing "
                 "price and its exponential moving average (EMA)."),
    build=build,
    default_params={"lookback_ema": 12, "short_lookback": 20},
    stream_type="pnl_daily",
    stream_category="d1",
    tags=["momentum", "ema", "d1"],
)
\end{codiceInterno}

\subsection{Approcci e overlay: lo stesso registro, due elenchi}

Il registro degli approcci serve \emph{due} cartelle, \pkgApprocci{} e
\pkgOverlay, e il package si deduce da dove viene la cartella:

\begin{codiceInterno}{python}{Il package dipende dalla provenienza}{approaches/registry.py:56}
    # il package dipende dalla cartella di provenienza: approaches/ oppure overlays/
    pkg = approach_dir.parent.name if approach_dir.parent.name in ("approaches", "overlays") else "approaches"
    module_dotted = f"{pkg}.{name}.approach"
\end{codiceInterno}

\noindent
Di conseguenza ci sono tre funzioni pubbliche con semantiche distinte:

\begin{center}
\begin{tabular}{@{}L{32mm}L{106mm}@{}}
\toprule
\textbf{Funzione} & \textbf{Che cosa restituisce, e perché} \\
\midrule
\rigaTenue
\cod{discover\_approaches()} & Solo \pkgApprocci. Sono quelli eseguibili dal
  pulsante \pulsante{Run backtest}. \\
\cod{discover\_overlays()} & Solo le cartelle \file{overlays/*/approach.py}.
  Non sono eseguibili da sole. \\
\rigaTenue
\cod{discover\_all()} & L'unione, \textbf{solo per risolvere un nome}, mai
  per costruire un elenco da mostrare. \\
\bottomrule
\end{tabular}
\end{center}

\begin{perche}{Perché gli elenchi sono separati}
Il razionale è scritto nel registro: mettere gli overlay fra i backtest
lanciabili «li faceva comparire fra i backtest lanciabili, dove sarebbero
falliti o --- peggio --- avrebbero prodotto uno stream vuoto senza dirlo». Uno
stream vuoto che non si lamenta è peggio di un errore, perché entra
nell'ensemble e lo cambia.
\end{perche}

\begin{nota}{Il campo \cod{family} non esiste nella dataclass}
\cod{\_scansiona} attacca all'oggetto un attributo \cod{family} che non è
dichiarato in \simbolo{ApproachDescriptor} e non compare in \cod{to\_dict()}:

\begin{codiceInterno}{python}{}{approaches/registry.py:113}
        # Da dove viene: e' l'unico modo affidabile per distinguere un
        # approccio da un overlay, perche' il contratto del file e' identico.
        try:
            desc.family = etichetta
        except Exception:      # dataclass congelata: non e' un problema
            pass
\end{codiceInterno}

\noindent
Funziona perché la dataclass non è congelata. Chi legge
\file{descriptor.py} non sa che esiste.
\end{nota}

\begin{pratica}{Aggiungere un approccio}
\begin{enumerate}
  \item Cartella senza \cod{\_}, con \file{approach.py} e
        \file{\_\_init\_\_.py}.
  \item \cod{def build(params: dict) -> pd.Series}: indice di date tz-naive,
        valori di rendimento giornaliero in frazione, causale.
  \item \cod{APPROACH\_DESCRIPTOR} con almeno \cod{name}, \cod{description},
        \cod{build}.
  \item Controllare che \cod{stream\_type} e \cod{stream\_category} stiano
        nelle due allowlist.
  \item Provarlo: \file{python -m execution.run\_approach <nome>}.
\end{enumerate}
\end{pratica}
